% 参考文献：条目照录原文，不译（见译校说明第五节第 5 条）
\renewcommand{\refname}{参考文献}
\addcontentsline{toc}{section}{参考文献}
\markboth{参考文献}{参考文献}

\begin{thebibliography}{[TT13]}

\bibitem[EH00]{EH00}
David Eisenbud and Joe Harris.
\newblock \emph{The geometry of schemes}.
\newblock Springer, 2000.

\bibitem[Gro66]{Gro66}
Alexander Grothendieck.
\newblock \emph{\'El\'ements de g\'eom\'etrie alg\'ebrique: IV. \'Etude locale des sch\'emas et des
  morphismes de sch\'emas, troisi\`eme partie}.
\newblock Publications Math\'ematiques de l'IH\'ES, 28:5--255, 1966.

\bibitem[Sta25]{Sta25}
The Stacks project authors.
\newblock \emph{The Stacks project}.
\newblock \url{https://stacks.math.columbia.edu}, 2025.

\bibitem[TT13]{TT13}
Robert W. Thomason and Thomas Trobaugh.
\newblock Higher algebraic $K$-theory of schemes and of derived categories.
\newblock In \emph{The Grothendieck Festschrift: a collection of articles written in honor of the
  60th birthday of Alexander Grothendieck}, pages 247--435. Springer, 2013.

\end{thebibliography}

\begin{yizhu}
四条文献在本书中的角色不同，读者可据此决定是否回溯：\textup{[EH00]} 是概形理论的入门教材，
例子与几何直观最多，与本书 §0、§6、§14 的动机部分最接近；\textup{[Gro66]} 即通称的 EGA IV$_3$，
是概形局部理论的原始文献，本书关于忠实平坦、形式纤维与丰沛性的若干论断可在此溯源；
\textup{[Sta25]} 是持续更新的在线百科，本书译注所引 tag 页码即出自它；\textup{[TT13]} 讨论概形的
高阶 $K$-理论与导出范畴，属附录 A 所铺垫的范畴论语言在更深层次上的应用，本书正文并未直接调用。
[依据：原书 p.171 文献表；各条性质据其标题与出版信息判定。]
\end{yizhu}
